\documentclass[11pt,a4paper,twoside,openany]{book}

\raggedright %izquierda
\usepackage[margin=2.54cm]{geometry}


%sangria
\usepackage{indentfirst}

\usepackage{parskip}
\setlength{\parskip}{0pt} % Sin espacio entre párrafos
\setlength{\parindent}{1.27cm} % Sangría de 1,27 cm para párrafos

% habilitamos Arial 11 pt en todo el documento (solo fontspec, sin mathptmx/t1enc)
\usepackage{anyfontsize}
\usepackage{fontspec}
\usepackage{url}

\usepackage{chngcntr}
\counterwithout{figure}{chapter}

\setmainfont{Arial}[
  Path           = ./fuentes/,
  UprightFont    = arial.ttf,
  BoldFont       = arialbd.ttf,
  ItalicFont     = ariali.ttf,
  BoldItalicFont = arialbi.ttf
]
\setsansfont{Arial}[
  Path           = ./fuentes/,
  UprightFont    = arial.ttf,
  BoldFont       = arialbd.ttf,
  ItalicFont     = ariali.ttf,
  BoldItalicFont = arialbi.ttf
]

%Habilitamos el TOC hasta subsubsection
\setcounter{tocdepth}{3}
\setcounter{secnumdepth}{3}

\newcommand\tab[1][1cm]{\hspace*{#1}}

\usepackage{setspace}
\doublespacing % Interlineado doble general; setspace deja notas al pie y flotantes (tablas) en sencillo

\usepackage{afterpage}
\usepackage{ragged2e}
\usepackage[none]{hyphenat}
\usepackage{xpatch}
\usepackage{blindtext}
%\usepackage{sectsty}
\usepackage{appendix}


%\def\UrlBreaks{\do\/\do-}
\usepackage{hyperref}
%\usepackage{breakurl}
%\usepackage[breaklinks]{hyperref} % Referencias e índices son enlaces al contenido. La opción hidelinks hace que se vean como texto normal


\usepackage{newclude} % Permite usar \include* que es igual que \include pero sin saltos de página adicionales
%\includeonly{cap01} % con este comando solo se incluirán los archivos mencionados haciendo que la compilación sea mucho más rápida. Referencias a otros capítulos quedarán rotas, se recomienda una compilación limpia cada vez que se cambie esta línea (es decir, borrar ficheros intermedios cuando se descomenta/comenta/cambia el argumento)

%\usepackage[top=3.0cm, bottom=2.0cm,left= 2.5cm,right=2cm]{geometry}

\usepackage{csquotes} % Paquete recomendado por biblatex cuando se usa babel

\usepackage{graphicx}
\usepackage[format=plain,justification=raggedright]{caption} % Justificación a la izquierda y formato plain

% Configuración para el formato de las figuras
\captionsetup[figure]{
  format=plain, % Estilo de formato
  labelformat=simple, % Formato de la etiqueta
  labelsep=none, % Sin separador entre la etiqueta y el título
  font=small, % Tamaño de la fuente
  justification=raggedright, % Justificación del texto a la izquierda
  singlelinecheck=false, % Permite justificación en varias líneas
  labelfont=bf, % Número de la figura en negrita
  textfont=it, % Título en cursiva
  skip=12pt, % Espacio entre la figura y la leyenda
  %parskip=10pt %es
}


% Personalización para que el título aparezca arriba de la figura
\usepackage{float} % Necesario para la opción [H]

\usepackage{amsmath,amssymb}
\usepackage{changepage}
\usepackage{subcaption}

\usepackage[table]{xcolor} % Necesaria la opción table para usar filas de colores en las tablas.
\usepackage{lipsum} % Para usar texto de relleno, sirve para tener una idea de cómo queda

\usepackage{booktabs} % para los comandos toprule, midrule, cmidrule y bottomrule: líneas horizontales de las tablas
\usepackage{threeparttable} % Entorno para las notas a pie de tabla

%--- Paquete algorithm y personalizaciones ---%
\usepackage{algorithm,algorithmic}
\floatname{algorithm}{Algoritmo}
\renewcommand{\listalgorithmname}{\'INDICE DE ALGORITMOS}

\usepackage{listings}
\renewcommand{\lstlistingname}{C\'odigo}% Listing -> Algorithm
\renewcommand{\lstlistlistingname}{\'INDICE DE CODIGOS DE PROGRAMA}

\usepackage{multirow}

\usepackage{tikz}

%--- Formato de títulos ---
\usepackage{titlesec} % Librería para cambiar el formato de \chapter, \section, ...
\titleformat{\chapter}[display]
{\bfseries\fontsize{11pt}{13.2pt}\selectfont}
{\centering\chaptertitlename \thechapter}{24pt}{}
\titlespacing*{\chapter}
{0pt}{30pt}{40pt}
\renewcommand{\chaptertitlename}{\fontsize{24pt}{24}\selectfont CAP\'ITULO }
\newpagestyle{mystyle}{\setfoot[][\thepage][]{}{\thepage}{}}
\pagestyle{mystyle}
\titleformat{\paragraph}[runin]{\normalfont\normalsize\itshape}{\theparagraph}{1.5em}{} % formato de \paragraph: fuente y tamaño normales y cursiva. 1em de separación
\titleformat{\subparagraph}[runin]{\normalfont\normalsize\itshape}{\thesubparagraph}{1.5em}{} % formato de \subparagraph: fuente y tamaño normales y cursiva. 1em de separación

%%%%%%%%%%%%%%%%%%%%%%%%%%%%% CONFIGURACIÓN DE LOS TÍTULOS %%%%%%%%%%%%%%%%%%%%%%%%%%%%%

% Configuración del formato de \section
\titleformat{\section}
  {\normalfont\bfseries} % Fuente normal, negrita y tamaño grande
  {\thesection}{1em} % Número de la sección
  {} % Espacio entre número y título

% Configuración del formato de \subsection
\titleformat{\subsection}
  {\normalfont\itshape\bfseries} % Fuente normal, cursiva, negrita y tamaño grande
  {\thesubsection}{1em} % Número de la subsección
  {} % Espacio entre número y título

% Configuración de \subsubsection como un título "run-in"
\titlespacing*{\subsubsection}
  {2em} % Sangría izquierda para el título
  {0pt} % Espacio antes del título
  {1em} % Espacio después del título

\titleformat{\subsubsection}[runin]
  {\normalfont\bfseries} % Formato del título
  {\thesubsubsection.} % Definir el número del título
  {0.5em} % Espacio entre el número y el título
  {} % Código antes del título
  
%Renombrar los epigrafos
\renewcommand{\tablename}{Tabla}
\renewcommand{\figurename}{Figura}

%--- AJUSTES MANUALES PARA CODIGOS ---%

\usepackage{listings}
\usepackage{color}

\definecolor{dkgreen}{rgb}{0,0.6,0}
\definecolor{gray}{rgb}{0.5,0.5,0.5}
\definecolor{mauve}{rgb}{0.58,0,0.82}

\lstset{frame=tb,
  language=Java,
  aboveskip=3mm,
  belowskip=3mm,
  showstringspaces=false,
  columns=flexible,
  basicstyle={\footnotesize\ttfamily\justifying},
  numbers=none,
  numberstyle=\tiny\color{gray},
  keywordstyle=\color{blue},
  commentstyle=\color{dkgreen},
  stringstyle=\color{mauve},
  breaklines=true,
  breakatwhitespace=true,
  tabsize=3
 }

% En el front-end al final de los usepackages
\usepackage[natbib=true,backend=biber,sorting=nyt,style=apa]{biblatex}
\renewcommand*{\bibfont}{\normalfont\normalsize} % Arial 11 pt, igual que el cuerpo
% Sangría francesa explícita de 1,27 cm en cada referencia
\setlength{\bibhang}{1.27cm}
\setlength{\bibitemsep}{0pt}
\AtBeginBibliography{%
  \setlength{\leftskip}{0pt}%
  \setlength{\parindent}{0pt}%
}

% add your references to this file
% Si pusieron otro nombre al archivo BiBteX usen ese :P
\addbibresource{referencias.bib}

\renewcommand\maketitle{}


%%%%%%%%%%%%%%%%%%%%%%%%%%%%% DOCUMENTO %%%%%%%%%%%%%%%%%%%%%%%%%%%%%
\begin{document}
\sloppy 
\maketitle

%%%%%%%%%%%%%%%%%%%%%%%%%%%%% TÍTULO, RESUMEN, AGRADECIMIENTOS Y DEDICATORIA %%%%%%%%%%%%%%%%%%%%%%%%%%%%%
%   TÍTULO
% Página con el título, hay que entrar en title.tex para modificar: nombre del trabajo, autor y fecha

\begin{titlepage}

\newcommand{\HRule}{\rule{\linewidth}{0.5mm}} % Defines a new command for the horizontal lines, change thickness here

\center % Center everything on the page
 
%	HEADING SECTIONS
%\includegraphics[width=2.25cm]{recursos/FACINFO}
%  \hspace{8cm}
%\includegraphics[width=2cm]{recursos/logoupm.png}
%\\[1cm]
%Nombre de la Universidad
{\fontsize{16}{20}\selectfont \textbf{Escuela Superior Polit\'ecnica del Litoral}}
\\[3.0em]
%Nombre de la Facultad
{\fontsize{16}{20}\selectfont \textbf{Facultad de Ingenier\'ia en Electricidad y Computaci\'on}}\\[3.0em]


%	TITLE SECTION
%TODO Ponerle título a la tesis
%% TÍTULO DEL PROYECTO INTEGRADO
{\fontsize{14}{18}\selectfont  Sistema de Generación de Perfiles del Personal Docente y Administrativo en ESPOL para la asignación inteligente de tareas}\\[3.0em]
{\fontsize{16}{20}\selectfont \textbf{Proyecto de Titulación}}\\[1.5em] 
{\fontsize{16}{20}\selectfont Previo la obtenci\'on del T\'itulo de:}\\[1.5em] 
{\fontsize{16}{20}\selectfont \textbf{Magíster en Ciencia de Datos}}\\[3.5em] 

 %	AUTHOR SECTION
{\fontsize{16}{20}\selectfont Presentado por:}\\[1.5em]
\author{Nombres1 y Apellidos1} {\fontsize{16}{20}\selectfont  Carlos Luis Sánchez Acosta}\\[1.5em]
% Eliminar o agregar esta sección si son más de dos autores
\author{Nombres2 y Apellidos2} {\fontsize{16}{20}\selectfont Pamela Nayeli Rugel D\'iaz}

\vspace{3.0em} %Esto hace la separación a la ciudad (NO BORRAR)
{\fontsize{16}{20}\selectfont Guayaquil - Ecuador}\\[1.5em]
{\fontsize{16}{20}\selectfont A\~no: 2026}

\thispagestyle{empty}
\end{titlepage}
%\let\clearpage\relax

% DEDICATORIA,AGRADECIMIENTO,DECLARACION_EXPRESA Y FIRMAS
% Título centrado con una línea debajo que cubre todo el ancho del texto
\titleformat{\chapter}[block]
  {\normalfont\normalsize\filcenter}
  {}{0pt}{}
  [\vspace{2pt}\titlerule]
%	DEDICATORIA
\chapter*{Dedicatoria}
\thispagestyle{empty}
\begin{flushright}
    \begin{minipage}{.45\linewidth}
    \raggedright
        \textit{Carlos Sánchez}\\
        Dedico este trabajo a 
        \\[3.0em]
        
        \textit{Pamela Rugel}\\
       Dedico este trabajo a 
        
    \end{minipage}
\end{flushright}

$\ $
\thispagestyle{empty}


%	AGRADECIMIENTOS
\chapter*{Agradecimientos}
\thispagestyle{empty}
\begin{flushright}
    \begin{minipage}{.45\linewidth}
    \raggedright
        \textit{Carlos Sánchez}\\
        Agradezco profundamente a 
        \\[3.0em]
        
        \textit{Pamela Rugel}\\
        Agradezco profundamente a 
    \end{minipage}
\end{flushright}

$\ $
\thispagestyle{empty}


%	DECLARACION EXPRESA
\chapter*{Declaración Expresa}
\thispagestyle{empty}

Nosotros Carlos Luis Sánchez Acosta y Pamela Nayeli Rugel Díaz acordamos y reconocemos que:
La titularidad de los derechos patrimoniales de autor (derechos de autor) del proyecto de graduación corresponderá al autor o autores, sin perjuicio de lo cual la ESPOL recibe en este acto una licencia gratuita de plazo indefinido para el uso no comercial y comercial de la obra con facultad de sublicenciar, incluyendo la autorización para su divulgación, así como para la creación y uso de obras derivadas. En el caso de usos comerciales se respetará el porcentaje de participación en beneficios que corresponda a favor del autor o autores. 

La titularidad total y exclusiva sobre los derechos patrimoniales de patente de invención, modelo de utilidad, diseño industrial, secreto industrial, software o información no divulgada que corresponda o pueda corresponder respecto de cualquier investigación, desarrollo tecnológico o invención realizada por mí/nosotros durante el desarrollo del proyecto de graduación, pertenecerán de forma total, exclusiva e indivisible a la ESPOL, sin perjuicio del porcentaje que nos corresponda de los beneficios económicos que la ESPOL reciba por la explotación de nuestra innovación, de ser el caso. 

En los casos donde la Oficina de Transferencia de Resultados de Investigación (OTRI) de la ESPOL comunique los autores que existe una innovación potencialmente patentable sobre los resultados del proyecto de graduación, no se realizará publicación o divulgación alguna, sin la autorización expresa y previa de la ESPOL. 

\vspace{1cm}

Guayaquil, 10 de Octubre del 2026.
\vspace{2.5cm}

\begin{table}[h!]
    \centering
    \begin{tabular}{ccc}
        % Si tienes las imágenes de las firmas, descomenta las siguientes líneas y ajusta la ruta:
        % \includegraphics[height=1.5cm]{recursos/firma_kevin.png} & \hspace{1cm} & \includegraphics[height=1.5cm]{recursos/firma_pamela.png} \\
        \rule{6cm}{0.5pt} & \hspace{1.5cm} & \rule{6cm}{0.5pt} \\
        Carlos Luis Sánchez Acosta & & Pamela Nayeli Rugel Díaz \\
    \end{tabular}
\end{table}
        
$\ $
\thispagestyle{empty}


%	EVALUADORES
\chapter*{Evaluadores}
\thispagestyle{empty}
\vspace{15.0em}
\begin{minipage}{.45\linewidth}
    \begin{flushleft} 
    \end{flushleft}
    \line(1,0){200} \\ \centering \textbf{Ing. Ronald Ra\'ul Criollo Bonilla} \\ PROFESOR DE LA MATERIA 
\end{minipage}
\hfill
\begin{minipage}{.45\linewidth}
    \begin{flushright}
    \end{flushright} 
    \line(1,0){200} \\ \centering \textbf{Ph.D. Federico Dom\'inguez Bonini} \\ PROFESOR TUTOR
\end{minipage}


$\ $
\thispagestyle{empty}
%\let\clearpage\relax

\let\cleardoublepage\clearpage
\frontmatter
\pagenumbering{Roman}

% RESUMEN Y ABSTRACT
% Resumen: solo centrado. Abstract: centrado y en cursiva
\newcommand{\formatoTituloResumen}[1]{\ifstrequal{#1}{Abstract}{\textit{#1}}{#1}}
\titleformat{\chapter}[block]
  {\normalfont\normalsize\filcenter}
  {}{0pt}{\formatoTituloResumen}
%	RESUMEN
\chapter{Resumen}

Estas instrucciones servirán de guía para la preparación de los trabajos que se presentarán como requisito en el proceso de graduación de la Materia Integradora de la Unidad Académica. El resumen deberá contener entre 150 a 200 palabras, incluyendo los siguientes cuatro componentes, descritos de forma concisa, comprensible y redactado en estilo impersonal (tercera persona): 1) se empieza con una breve introducción, objetivos, hipótesis y justificación del proyecto descrito en tiempo presente; 2) desarrollo del proyecto, donde se describirán brevemente los materiales, equipos, técnicas, normas etc. utilizadas en el proyecto. Esta sección se redacta en tiempo pasado; 3) continúe con los resultados donde se describen de forma concisa solo los principales y más novedosos resultados escritos en tiempo pasado; 4) finalmente, se presentan las conclusiones generales del proyecto en tiempo presente. Además, deberá incluir al menos 4 palabras clave al final del documento. Todo el resumen se presentará en un sólo cuerpo. Utilice el contador de palabras del procesador de texto para asegurarse del tamaño del documento. 

\textbf{Palabras Clave:} Formato, Proyecto Integrador, etc. (mínimo 4 y máximo 5 palabras).
Procure que las palabras claves que aquí coloque no repitan las que ya se encuentran incluidas en su título.


%	ABSTRACT
\chapter{Abstract}
These instructions will serve as a guide for preparing the papers that will be submitted as a requirement for graduation in the Integrative Subject of the Academic Unit. The abstract should contain between 150 and 200 words, including the following four components, described concisely, comprehensibly, and written in an impersonal style (third person): 1) Begin with a brief introduction, objectives, hypothesis, and justification of the project, described in the present tense; 2) Project development, where the materials, equipment, techniques, standards, etc., of the project will be briefly described. This section is written in the past tense; 3) Continue with the results, where only the main and most novel results are concisely described, written in the past tense; 4) Finally, present the general conclusions of the project in the present tense. In addition, include at least four keywords at the end of the document. The entire abstract should be presented in a single document. Use the word counter in your word processor to ensure the document size is correct.

\textbf{Keywords:} Gesture Detection, Automatic Feedback, Video Processing, Artificial Intelligence, Data Analysis.

\newpage
%\let\clearpage\relax

% Se restablece el formato anterior de \chapter para los índices
\titleformat{\chapter}[display]
{\bfseries\fontsize{11pt}{13.2pt}\selectfont}
{\centering\chaptertitlename \thechapter}{24pt}{}

%   ÍNDICE
%  INDICE GENERAL FORMATO
\renewcommand*{\contentsname}{\'Indice General}
\phantomsection
\addcontentsline{toc}{chapter}{\contentsname}
\tableofcontents  % indice de contenidos

% ABREVIATURAS
\chapter{Abreviaturas}
\renewcommand{\arraystretch}{2}
\begin{table}[htbp]
\begin{tabular}{ll}
BiLSTM          & Long Short-Term Memory Bidireccional \\
BiLSTM-CNN      & Long Short-Term Memory Bidireccional with Convolutional Neural Network \\
CNN             & Convolutional Neural Network\\
Conv-LSTM       & Convolutional Long Short-Term Memory\\
CPM             & Capacidad, Planificación y Monitoreo\\
CTI             & Centro de Tecnolog\'ia de Informaci\'on\\
DB-LSTM         & Densely-Connected Bi-directional Long Short-Term Memory\\
ESPOL           & Escuela Superior Polit\'ecnica del Litoral \\
FC              & Fully Connected \\
HAR             & Human Action Recognition\\
LSTM            & Long Short-Term Memory\\
MMT             & Multimodal Transformer Model\\
NTU RGB+D       & Dataset RGB+D de la Universidad Tecnológica de Nanyang \\
\end{tabular}
\end{table}

\newpage

%   INDICE DE FIGURAS, TABLAS Y ALGORITMOS
\renewcommand*{\listfigurename}{\'Indice de figuras}
\renewcommand{\thefigure}{\arabic{figure}}
\counterwithout{figure}{chapter}
\phantomsection
\addcontentsline{toc}{chapter}{\listfigurename}
\listoffigures

%\renewcommand*{\listtablename}{\'INDICE DE TABLAS}
%\addcontentsline{toc}{chapter}{\listtablename}
%\listoftables

% Si no hay muchos algoritmos no tiene mucho sentido esta lista
%\addcontentsline{toc}{chapter}{\listalgorithmname}
%\listofalgorithms 

%\addcontentsline{toc}{chapter}{\lstlistlistingname}
%\lstlistoflistings 

%%%%%%%%%%%%%%%% CUERPO DE LA PÁGINA DEL CAPITULO %%%%%%%%%%%%%%%%%%
\mainmatter

% Configuración del formato del capítulo
\titleformat{\chapter} % Formato colgante
  {\bfseries\fontsize{11pt}{13.2pt}\selectfont} % Formato de fuente 
  {\thechapter.}% Mostrar solo el número del capítulo seguido de un punto
  {1em}% Espacio entre el número y el título
  {}% Código antes del título

  
% TABLAS, FIGURAS, EXPRESIONES MATEMÁTICAS Y ALGORITMOS
\newpage
\thispagestyle{empty}
\vspace*{6cm}
\begin{center}
    \fontsize{16pt}{19.2pt}\selectfont
    \textbf{Cap\'itulo 1}
\end{center}
\vspace*{\fill}
\phantomsection
\addcontentsline{toc}{chapter}{Cap\'itulo 1}

\include*{capitulos/cap01}
%\let\clearpage\relax

% CONTENIDOS DEL TFG
\newpage
\thispagestyle{empty}
\vspace*{6cm}
\begin{center}
    \fontsize{16pt}{19.2pt}\selectfont
    \textbf{Cap\'itulo 2}
\end{center}
\vspace*{\fill}
\phantomsection
\addcontentsline{toc}{chapter}{Cap\'itulo 2}

\include*{capitulos/cap02}
%\let\clearpage\relax

% REFERENCES
\newpage
\thispagestyle{empty}
\vspace*{6cm}
\begin{center}
    \fontsize{16pt}{19.2pt}\selectfont
    \textbf{Cap\'itulo 3}
\end{center}
\vspace*{\fill}
\phantomsection
\addcontentsline{toc}{chapter}{Cap\'itulo 3}

%\include*{capitulos/cap03}
%\let\clearpage\relax

% CONCLUSIONES Y LÍNEAS FUTURAS DE TRABAJO
\newpage
\thispagestyle{empty}
\vspace*{6cm}
\begin{center}
    \fontsize{16pt}{19.2pt}\selectfont
    \textbf{Cap\'itulo 4}
\end{center}
\vspace*{\fill}
\phantomsection
\addcontentsline{toc}{chapter}{Cap\'itulo 4}

%\include*{capitulos/cap04}
%\let\clearpage\relax
% AGREGAR MAS CAPITULO
%\include{capitulos/cap05}


%%%%%%%%%%%%%%%%%%%%%%%%%%%%% BACKMATTER %%%%%%%%%%%%%%%%%%%%%%%%%%%%%
\backmatter  
\titleformat{\chapter}[display]
{\bfseries\centering\normalsize} 
{\centering\chaptertitlename \thechapter}{24pt}{}
\titlespacing*{\chapter}
{0pt}{30pt}{40pt}
%%% Bibliografía bibname
% Referencias y Apéndices sin número de página visible (el contador no se toca)
\clearpage
\pagestyle{empty}
\assignpagestyle{\chapter}{empty}
\renewcommand{\bibname}{Referencias}
\phantomsection
\addcontentsline{toc}{chapter}{Referencias} % Para que salga en el índice general
%\bibliographystyle{ieeetr}
%\bibliography{referencias}

\printbibliography{}

%%%%%%%%%%%%%%%%%%%%%%%%%%%%% APÉNDICES %%%%%%%%%%%%%%%%%%%%%%%%%%%%%

%\appendix
%\clearpage
%\renewcommand{\appendixtocname}{Ap\'endices}
%\renewcommand{\appendixname}{Ap\'endices}
%\renewcommand{\appendixpagename}{Ap\'endices}
%\renewcommand{\thesection}{\Alph{section}}
%\addappheadtotoc
%\appendixpage
%\pagestyle{empty}

%\include*{apendices}


\end{document}